% !TEX root = ../../../main.tex
% 所属：第一部分：单变量分析学 / 第1章：实数理论——从Peano公理到实数的连续性
% 节：1.5 $\Q$有理数的构建
% 正文按原文件顺序保存；由 ../chapter1.tex 调用。
\section{$\Q$有理数的构建}
\begin{construction}有理数的构建
	$$
	S=\left\{(a, b): a, b \in \mathbb{Z}\text{且}b\ne0\right\}
	$$
	
	
	在 $S$ 上定义 $\sim$ 关系
	\[(a, b) \sim(c, d) \text { iff } a d=b c\]
	\[\Q:=S/\sim\]
	
\end{construction}



现验证 $\sim$ 是等价关系
\begin{proof}
	\begin{enumerate}
		\item 自反性 $(a, b) \sim(a, b) \Leftrightarrow a b=b a$
		\item   对称性 $(a, b) \sim ( c, d) \Leftrightarrow a d=b c \Leftrightarrow c b=d a$
		\[\text{故}(c, d) \sim(a, b)\]
		
		\item  传递性 $(a, b) \sim(c, d)$ 且 $(c, d) \sim(e, f)$
		
		$$
		\begin{aligned}
			& \Leftrightarrow a d=b c \text { 且 } c f=d e \\
			& \Leftrightarrow a f=b e \Leftrightarrow(a, b) \sim(e, f)
		\end{aligned}
		$$
		
	\end{enumerate}
\end{proof}
每个等价类$[(a, b)]$ 是一个有理 数.

\begin{definition}$\Q$上的运算\\
	加法 $[(a, b)]+[(c, d)]:=[( a d+b c, b d)]$\\
	加法逆元 $-[(a, b)]:=[(-a, b)]$ 或$[(a,-b)]$\\
	乘法 $[(a, b)] \cdot[c, d]:=[(a c, b d)]$\\
	乘法逆元 $[(a, b)]^{-1}=[(b, a)]$,其中$a,b\ne0$
\end{definition}

在$\Q$上$[(a, b)]$ 记作 $\dfrac{a}{b} \quad$


显然 $[(n, 1)]=n \in \mathbb{Z} \Rightarrow \mathbb{Z} \subset\Q$

$\Q$是一个数域

\begin{example}
	$$\dfrac{1}{2}=[(2,4)]=[(3,6)]=\cdots$$
	$$\dfrac{3}{5}=[(6,10)]=[(3,5)]=\cdots$$
\end{example}

